\chapter{原点、空间群与模型约束}
\section{对称关系必须同时作用于相位}
设分数坐标对称操作为 $x\mapsto Rx+t$。保持晶体密度不变意味着
\begin{equation}
 F(hR)=\exp(-2\pi\iu h\cdot t)F(h),
\end{equation}
此处 $h$ 采用代码中的行向量记号。若 P1 候选相对于标准设置还差一个原点平移 $o$，其结构因子会乘上 $e^{2\pi\iu h\cdot o}$，于是应检验
\begin{equation}\label{eq:sg-phase}
 F(hR)\exp\{2\pi\iu[h\cdot t+(hR-h)\cdot o]\}-F(h)=0.
\end{equation}
这个公式同时解释了三件事：空间群判断需要相位；原点有时必须搜索；纯平移导致的消光不能通过移动原点消除。

\section{搜索的是候选假设}
新模块对显式给出的候选空间群做度量兼容性检查，然后最小化式\eqref{eq:sg-phase}的归一化复残差。原点的 Jacobian 直接来自指数函数：
\begin{equation}
 \frac{\partial r_{h,o}}{\partial o_j}
 =2\pi\iu(hR-h)_jF(hR)
 e^{2\pi\iu[h\cdot t+(hR-h)\cdot o]}.
\end{equation}
有限网格和由原子对得到的起点共同初始化局部优化。若 $hR-h$ 只张成二维空间，那么第三个原点方向没有被该群的相位关系固定；程序报告这个秩，而不假装找到了唯一原点。

P1 的相位不一致度恒为零，因此不能简单把所有群按这个数排序，再宣布最小值就是正确空间群。低不一致度只允许产生一个新的结构候选。最终还需要看系统消光、等价反射一致性、几何投影和独立选择数据。

\section{从 P1 全晶胞到非对称单元}
原点对齐之后，每个对称操作应把一组原子映到同元素的另一组原子。代码使用周期晶胞度量计算距离，并通过一对一分配避免多个原子同时认领同一个位置。若操作后的最大匹配距离超过声明的容差，该候选被拒绝。

将一个对称轨道内匹配的原子逆变换回代表位置，可得若干坐标估计：
\begin{equation}
 x_{jo}^{\rm back}=R_o^{-1}(x_{\pi_o(j)}+n_{jo}-t_o).
\end{equation}
整数平移 $n_{jo}$ 用最小周期距离确定，平均后得到对称化位置。位移张量也必须作相同的逆旋转。完成投影后，再展开全部物理位置，核对元素数、占有率和坐标偏差。

对斜晶胞，直接把每个分数坐标差四舍五入未必得到最近的晶格平移。保留的 \code{Cell.minimum_image} 用三维球形分支限界求最近晶格向量，供配对、几何限制和验证统一使用。

\begin{worked}{任意原点下的反演结构}
将一个斜晶胞中的 $P\overline1$ 六原子全晶胞任意平移后，用原点搜索和对称投影可以恢复三个原子的非对称单元。本次独立测试的相位不一致度约为 $1.95\times10^{-30}$，含原点相位因子后的结构因子差小于 $5.3\times10^{-14}$。这是恒等关系与实现一致性的检查，不是未知实验数据空间群识别率。
\end{worked}

\section{精确约束与柔性限制}
空间群规定的特殊位置是精确约束。例如反演中心上的坐标不应靠一个很强的弹簧“拉近”，而应从自由变量中消去。用 $x=x_0+Tv$ 表示允许的坐标切空间，可避免无意义的自由度和奇异法方程。

键长、键角、平面性通常属于模型知识的柔性限制。例如
\begin{equation}
 L_{\rm bond}=\frac12\left(\frac{d_{ij}(\theta)-d_0}{\sigma_d}\right)^2.
\end{equation}
$\sigma_d$ 控制允许偏离的程度。理想 MOF 的连接关系可以帮助选择 $d_0$，但它仍是先验知识。最终报告应分开给出衍射残差和限制残差，并检查某个参数是否主要由限制而不是数据确定。
